\documentclass[14pt,a4paper]{extarticle}
\usepackage[T2A]{fontenc}
\usepackage[utf8]{inputenc}
\usepackage[english,russian]{babel}
\usepackage{lmodern}
\usepackage{geometry}
\usepackage{array}
\usepackage{tabularx}
\usepackage{longtable}
\usepackage{ltablex}
\usepackage{booktabs}
\usepackage{hyperref}
\usepackage{indentfirst}
\usepackage{microtype}
\usepackage{graphicx}
\usepackage{setspace}
\usepackage{listings}
\usepackage{graphicx}
\usepackage{caption}
\usepackage{float}
\usepackage{tocloft}
\usepackage{titlesec}
\usepackage{newtxtext,newtxmath}

\titleformat{\section}        % форматируем \section
  {\normalfont\bfseries}      % стиль: обычный + жирный
  {\thesection.}              % номер секции с точкой
  {2mm}                       % расстояние между номером и текстом
  {}                           % перед заголовком можно добавить что-то
\titlespacing*{\section}{1.5cm}{*3}{*1}  % {отступ слева}{отступ сверху}{отступ снизу}

\titleformat{\subsection}        % форматируем \section
  {\normalfont\bfseries}      % стиль: обычный + жирный
  {\thesubsection.}              % номер секции с точкой
  {2mm}                       % расстояние между номером и текстом
  {}                           % перед заголовком можно добавить что-то
\titlespacing*{\subsection}{1.5cm}{*3}{*1}  % {отступ слева}{отступ сверху}{отступ снизу}

\captionsetup{justification=centering}
\renewcommand{\ttdefault}{cmtt}
\keepXColumns
\setlength{\parindent}{3em}
\newgeometry{left=2.5cm,right=2.5cm,top=3cm,bottom=2.5cm}

\renewcommand{\contentsname}{\hspace*{4em}{\smallСодержание}}
\setlength{\cftbeforesecskip}{0.5em}
\renewcommand{\cftsecleader}{\cftdotfill{\cftdotsep}}
\renewcommand{\cftsecaftersnum}{.\hspace{2mm}}

\title{Отчёт по "Управление ИТ-проектами"}
\author{Выполнил студент группы М-РИТ-25-1, Станиславчук Сергей}

\begin{document}

\setstretch{1.5}
\noindent
\begin{minipage}[t]{0.01\textwidth}
    \includegraphics{~/programmer/++Programmer/5_1/LSTU.jpg}
\end{minipage}%
\begin{minipage}[c]{1.1\linewidth}
    \centering
    {\small
    МИНИСТЕРСТВО НАУКИ И ВЫСШЕГО ОБРАЗОВАНИЯ РФ \\[1em]
    ФЕДЕРАЛЬНОЕ ГОСУДАРСТВЕННОЕ БЮДЖЕТНОЕ \\
    ОБРАЗОВАТЕЛЬНОЕ УЧРЕЖДЕНИЕ ВЫСШЕГО ОБРАЗОВАНИЯ \\
    «ЛИПЕЦКИЙ ГОСУДАРСТВЕННЫЙ ТЕХНИЧЕСКИЙ УНИВЕРСИТЕТ»
    }
\end{minipage}%

\vspace{3em}

\noindent
\hspace{1em}%
\begin{minipage}[t]{0.8\textwidth}
\textbf{Институт} \hspace{2em} компьютерных наук \\
\textbf{Кафедра} \hspace{2em} автоматизированных систем управления
\end{minipage}


\vspace{6em}

\begin{center}
\textbf{Практическая работа №11} \\
по дисциплине \\ <<Управление ИТ-проектами>>
\end{center}

\vspace{6em}

\noindent
\begin{tabular*}{\textwidth}{@{\extracolsep{\fill}} l c r }
    Студент М-РИТ-25-1 & \makebox[4cm]{\hrulefill} & Станиславчук С.М. \\[2em]
Руководитель & \\
К.т.н., доцент & \makebox[4cm]{\hrulefill} & Ведищев В.В. \\
\end{tabular*}

\vfill
\begin{center}
    Липецк, 2026
\end{center}

\thispagestyle{empty}

\newpage
\tableofcontents
\newpage

\section{Управление разработкой через инструмент «тикетов»}

\subsection{Понятие «тикет»}

Тикет представляет собой единицу работы в системе управления задачами, фиксирующую конкретное требование, задачу или проблему. В рамках разработки системы корпоративного обмена сообщениями тикеты используются для формализации всех видов работ: реализации функциональности, исправления ошибок, доработки интерфейса и обеспечения безопасности.

Каждый тикет содержит описание задачи, приоритет, статус выполнения, ответственного исполнителя и при необходимости дополнительные материалы (логи, скриншоты, ссылки на код). Такой подход позволяет структурировать процесс разработки и обеспечить прозрачность выполнения работ.

\subsection{Понятие «характеристики взаимодействия»}

Характеристики взаимодействия определяют правила и параметры работы с тикетами внутри команды. К ним относятся порядок постановки задач, распределение ответственности, приоритеты выполнения и способы контроля.

В рамках проекта тикеты проходят через определённые статусы (например, «создан», «в работе», «на проверке», «завершён»), что позволяет отслеживать текущее состояние разработки. Также фиксируется связь между тикетами и изменениями в системе контроля версий, что обеспечивает прослеживаемость выполненных работ.

\subsection{Этапы управления разработкой}

Управление разработкой через тикеты осуществляется последовательно. На первом этапе формируется тикет с описанием задачи на основе требований. Далее задача назначается исполнителю, который выполняет реализацию в отдельной ветке системы контроля версий.

После завершения разработки тикет переводится на этап проверки, где осуществляется код-ревью и тестирование. В случае выявления ошибок тикет возвращается на доработку. После успешной проверки задача считается завершённой и интегрируется в основную ветку проекта.

Таким образом, тикеты выступают связующим элементом между требованиями, реализацией и тестированием.

\subsection{Поддержка итеративности этапов разработки}

Использование тикетов обеспечивает итеративный характер разработки. Проект разбивается на небольшие задачи, каждая из которых проходит полный цикл — от постановки до завершения.

Это позволяет постепенно наращивать функциональность системы, регулярно получать промежуточные результаты и оперативно реагировать на изменения требований. Итеративный подход также упрощает выявление и исправление ошибок, так как каждая задача рассматривается отдельно и имеет ограниченный объём.

В результате управление через тикеты повышает управляемость проекта, улучшает взаимодействие внутри команды и способствует более эффективной реализации системы.

\end{document}
